\section*{CHƯƠNG 1. GIỚI THIỆU}
\setcounter{section}{1}
\setcounter{subsection}{0}
\addcontentsline{toc}{section}{Chương 1. Giới thiệu}

\subsection{Bối cảnh và Đặt vấn đề}
Bài toán nhận dạng chữ số viết tay (Handwritten Digit Recognition) là một nhánh kinh điển trong lĩnh vực Nhận dạng mẫu và Thị giác máy tính. Đây là bài toán phân loại đa lớp có giám sát ($K = 10$ lớp, từ $0$ đến $9$). Mỗi mẫu ảnh mức xám hai chiều kích thước $28 \times 28$ pixel tương đương vector đặc trưng $\mathbf{x} \in \mathbb{R}^{784}$ với cường độ điểm ảnh chuẩn hóa về đoạn $[0, 1]$. Mục tiêu là tối ưu hóa hàm quyết định $f: \mathbb{R}^{784} \rightarrow \{0, \dots, 9\}$ nhằm cực tiểu hóa sai số phân loại trên tập kiểm thử độc lập.

Dữ liệu chữ số viết tay đặt ra ba thách thức kỹ thuật cốt lõi:
\begin{itemize}
    \item \textbf{Biến thiên hình thái học:} Phong cách viết tay biến thiên đa dạng về nét bút, góc nghiêng và độ co giãn.
    \item \textbf{Tương quan không gian 2D:} Các điểm ảnh lân cận liên kết chặt chẽ tạo thành nét vẽ hình học; việc duỗi phẳng ảnh thành vector $1\text{D}$ làm mất hoàn toàn cấu trúc không gian này.
    \item \textbf{Sự tương đồng giữa các lớp:} Nhiều cặp chữ số có nét vẽ tương đồng dễ gây nhầm lẫn (ví dụ: $4$ và $9$, $3$ và $5$, $7$ và $1$).
\end{itemize}

\subsection{Mục tiêu nghiên cứu và Phạm vi đề tài}
Mục tiêu trung tâm của đồ án là thực nghiệm toàn diện và đối chứng khách quan giữa 4 trường phái thuật toán trên bộ dữ liệu chuẩn MNIST:
\begin{enumerate}
    \item \textbf{Mô hình xác suất Bayesian (\code{1\_Bayes}):} Naïve Bayes đa thức (\code{MultinomialNB-NP}) với giả định độc lập có điều kiện và kỹ thuật làm trơn Laplace.
    \item \textbf{Mô hình tuyến tính cổ điển (\code{2\_Logistic}):} Hồi quy Logistic đa lớp (\code{LogisticRegression-NP}) theo chiến lược One-Vs-Rest kèm điều chuẩn $L_2$.
    \item \textbf{Mô hình Học sâu (\code{3\_MLP\_CNN}):} Mạng nơ-ron đa tầng (\code{MLP}) với lan truyền thuận/ngược tự cài và Mạng tích chập (\code{CNN}) khai thác không gian 2D.
    \item \textbf{Mô hình Học kết hợp (\code{4\_Ensemble}):} Thuật toán Boosting (\code{AdaBoost}) kết hợp tuần tự các bộ phân loại yếu với nhiều dạng kiến trúc cơ sở.
\end{enumerate}
Các thuật toán cốt lõi (Naïve Bayes, Logistic Regression, Backpropagation của MLP) được tự cài đặt từ đầu (\textit{from-scratch}) bằng NumPy, sau đó đối chứng định lượng với scikit-learn và PyTorch trên cùng một phân chia dữ liệu cố định ($50.000$ train, $10.000$ val, $10.000$ test).

\subsection{Cấu trúc báo cáo}
Báo cáo được tổ chức thành 7 chương: Chương 1 giới thiệu bối cảnh; Chương 2 mô tả dữ liệu và tiền xử lý; Chương 3 trình bày cơ sở lý thuyết và mã giả; Chương 4 phân tích siêu tham số và tính công bằng; Chương 5 trình bày kết quả thực nghiệm chi tiết và bảng tổng hợp; Chương 6 thảo luận chuyên sâu, đánh giá biên Pareto và 4 kết luận đối chiếu; Chương 7 kết luận và định hướng phát triển.
